\chapter{Clebsch--Gordan 系数}
\label{app:cg}

\section{怎样读系数表}
固定 $j_1,j_2$。表的每一行用两个角动量沿同一 $z$ 轴的投影标记，代表一个乘积基矢，
\[
    \ket{m_1,m_2} = \ket{j_1,m_1}_1 \otimes \ket{j_2,m_2}_2 .
\]
若制备的是总角动量态 $\ket{j_1,j_2;j,m}$，测出这对投影的振幅就是\dfnref{dfn:qm-cg}中的系数。相应的展开为
\begin{equation}
    \ket{j_1,j_2;j,m} = \sum_{m_1+m_2=m} C^{jm}_{j_1m_1,j_2m_2}\ket{m_1,m_2}.
\end{equation}
总投影 $m$ 决定哪些乘积态可以出现。表中每一列给出一个归一化耦合态，各项的模平方就是分别测量时的概率。
采用 Condon--Shortley 相位约定：在每个 $j$ 的最高态 $m = j$ 中，取允许的最大 $m_1$ 所对应的系数为正，随后用总降算符依次确定其余态的相位。

\paragraph*{两个自旋 $1/2$}
在 $m = 1$ 和 $m = -1$ 时分别只有 $\ket{\up,\up}$ 和 $\ket{\dn,\dn}$，对应的系数都为 $1$。
在 $m = 0$ 时，\exmref{exm:qm-two-spins}给出的系数为
\begin{equation}
    \begin{array}{c|cc}
        & j = 1 & j = 0 \\ \hline
        \ket{\up,\dn} & 1/\sqrt2 & 1/\sqrt2 \\
        \ket{\dn,\up} & 1/\sqrt2 & -1/\sqrt2
    \end{array}
\end{equation}
两列都预言两个出口各占一半。第二列的相对负号使它与第一列正交，对应总角动量为零的单态。

\paragraph*{轨道角动量 $\ell$ 与自旋 $1/2$}
左因子取轨道，右因子取自旋。固定总投影 $m$ 后，自旋向上时轨道投影为 $m - 1/2$，自旋向下时为 $m + 1/2$。
令
\begin{equation}
    a_m = \sqrt{\frac{\ell + m + 1/2}{2\ell + 1}},\qquad
    b_m = \sqrt{\frac{\ell - m + 1/2}{2\ell + 1}}.
\end{equation}
两个 $j$ 都允许时，系数为
\begin{equation}
    \begin{array}{c|cc}
        & j = \ell + 1/2 & j = \ell - 1/2 \\ \hline
        \ket{m - 1/2,\up} & a_m & -b_m \\
        \ket{m + 1/2,\dn} & b_m & a_m
    \end{array}
\end{equation}
在端点 $m = \pm(\ell + 1/2)$ 只取第一列和允许的乘积态，其系数为 $1$。$\ell = 0$ 时只有 $j = 1/2$。
例如，制备 $\mathrm{p}$ 电子（$\ell = 1$）于 $j = 3/2,m = 1/2$，读第一列得到
\begin{equation}
    \ket{1,1/2;3/2,1/2}
    = \sqrt{2/3}\ket{0,\up} + \sqrt{1/3}\ket{1,\dn}.
\end{equation}
因此测得自旋向上的概率为 $2/3$，此时轨道投影必为零。

\section{从最高态构造整张表}
\label{sec:cg-lowering}
先把乘积态按总投影 $m = m_1 + m_2$ 分组。不同组之间的系数为零，同一组里只需处理一张有限维的换基矩阵。
允许的 $j$ 从 $j_1 + j_2$ 逐个减一，直到 $\abs{j_1 - j_2}$。此外必须有 $\abs m \le j$。
即使这些条件都满足，系数也可能因对称性而为零，下面的例子会出现这样的零。

最高投影 $m = j_1 + j_2$ 要求两个投影分别取到最大值，对应的乘积态为 $\ket{j_1,j_1}_1 \otimes \ket{j_2,j_2}_2$。从这个态往下降，总角动量 $j$ 保持不变。
降算符作用在耦合态上只有一项，在乘积态上却有两项。令 $d(j,m) = \sqrt{(j + m)(j - m + 1)}$，则
\begin{equation}
    \begin{alignedat}{1}
        \hbar^{-1}\op J_-\ket{j_1,j_2;j,m}
            &= d(j,m)\ket{j_1,j_2;j,m - 1},\\
        \hbar^{-1}\op J_-\ket{m_1,m_2}
            &= d(j_1,m_1)\ket{m_1 - 1,m_2} + d(j_2,m_2)\ket{m_1,m_2 - 1}.
    \end{alignedat}
\end{equation}
每降一次，先把落到同一个乘积态的振幅相加，再除以 $d(j,m)$，便得到下一组系数。降到 $m = -j$ 时停止。
寻找下一个 $j$ 的最高态时，在 $m = j$ 的乘积态中，取与所有已经构造出的同投影态正交的归一化组合。
由\thmref{thm:qm-angular-addition}，此时恰好剩一个方向。用本附录开头的相位约定确定符号，就可以继续下降。

\paragraph*{两个角动量 $1$}
本例固定 $j_1 = j_2 = 1$，沿用正文的简记 $\ket{j,m}_J$ 表示耦合态，$\ket{m_1,m_2}$ 表示乘积态。例如 $\ket{2,2}_J = \ket{1,1}$。
在 $m = 1$ 组中，下降一次得到对称组合，和它正交的另一列就是 $j = 1$ 的最高态，
\begin{equation}
    \ket{2,1}_J = \frac{\ket{1,0} + \ket{0,1}}{\sqrt2},\qquad
    \ket{1,1}_J = \frac{\ket{1,0} - \ket{0,1}}{\sqrt2}.
\end{equation}
再降一次，$\ket{1,0}$ 与 $\ket{0,1}$ 都能到达 $\ket{0,0}$。两条路径在 $j = 2$ 列相加，在 $j = 1$ 列相消。
因此 $m = 0$ 的前三个乘积态按下表组合；最后一列由前两列的正交补给出。
\begin{equation}
    \begin{array}{c|ccc}
      m = 0 & j = 2 & j = 1 & j = 0 \\ \hline
      \ket{1,-1} & 1/\sqrt6 & 1/\sqrt2 & 1/\sqrt3 \\
      \ket{0,0}  & 2/\sqrt6 & 0 & -1/\sqrt3 \\
      \ket{-1,1} & 1/\sqrt6 & -1/\sqrt2 & 1/\sqrt3
    \end{array}
    \label{eq:cg-one-one}
\end{equation}
例如两边都制备为投影零，测总角动量时沿第二行横读，$j = 2,1,0$ 的概率依次为 $2/3,0,1/3$。
要补出 $m = -1,-2$ 的各列，可以继续下降，也可以直接用下一节的投影反号公式。

\section{逆变换、正交性与对称性}
固定 $m$，用行标 $r = (m_1,m_2)$、列标 $j$ 将系数简写为 $C_{rj}$。因为两组基都正交归一，换基矩阵是幺正矩阵。
本附录的相位约定使元素都是实数，因此逆变换就是转置。读一列给出耦合态，读一行给出乘积态，
\begin{subequations}
    \begin{align}
        \ket{j_1,j_2;j,m} &= \sum_r C_{rj}\ket r,\\
        \ket r &= \sum_j C_{rj}\ket{j_1,j_2;j,m}.
    \end{align}
\end{subequations}
分别对两式取内积便得到两种核对方法，同列或同行的平方和为一，不同列或不同行的内积为零，
\begin{equation}
    \sum_r C_{rj}C_{rj'} = \delta_{jj'},\qquad
    \sum_j C_{rj}C_{r'j} = \delta_{rr'}.
\end{equation}
交换两个因子的次序，或者同时把三个投影反号，系数分别满足
\begin{equation}
    \begin{alignedat}{1}
      C^{jm}_{j_2m_2,j_1m_1} &= (-1)^{j_1+j_2-j}C^{jm}_{j_1m_1,j_2m_2},\\
      C^{j,-m}_{j_1,-m_1,j_2,-m_2} &= (-1)^{j_1+j_2-j}C^{jm}_{j_1m_1,j_2m_2}.
    \end{alignedat}
\end{equation}
例如在式\eqref{eq:cg-one-one}中交换两行所标的投影，$j = 1$ 的系数变号。乘积态 $\ket{0,0}$ 在交换下回到自身，这一行的 $j = 1$ 系数只可能为零。

\sectionstar{任意两个角动量的显式公式}
只求一个指定系数时，可以直接用有限和。先检查三角条件、投影范围、各投影与相应角动量相差整数以及 $m = m_1 + m_2$。
还需 $j_1 + j_2 - j$ 为整数。违反任一条件便记系数为零。条件满足时，先算两个非负因子
\begin{equation}
    \begin{alignedat}{1}
        A &= \frac{(2j + 1)(j_1 + j_2 - j)!\,(j + j_1 - j_2)!\,(j - j_1 + j_2)!}{(j_1 + j_2 + j + 1)!},\\
        B &= (j + m)!\,(j - m)!\,(j_1 + m_1)!\,(j_1 - m_1)!\,(j_2 + m_2)!\,(j_2 - m_2)! .
    \end{alignedat}
\end{equation}
每一项的分母由六个阶乘组成，系数为
\begin{equation}
    \begin{alignedat}{1}
        D_k &= k!\,(j_1 + j_2 - j - k)!\,(j_1 - m_1 - k)!\,(j_2 + m_2 - k)!\\
           &\qquad{} \times (j - j_2 + m_1 + k)!\,(j - j_1 - m_2 + k)!,\\
        C^{jm}_{j_1m_1,j_2m_2} &= \sqrt{AB}\sum_{k=k_{\min}}^{k_{\max}}\frac{(-1)^k}{D_k}.
    \end{alignedat}
\end{equation}
允许的整数 $k$ 恰好使这些阶乘的自变量都非负，因此上下限可直接写成
\begin{equation}
    \begin{alignedat}{1}
        k_{\min} &= \max(0,\,j_2 - j - m_1,\,j_1 + m_2 - j),\\
        k_{\max} &= \min(j_1 + j_2 - j,\,j_1 - m_1,\,j_2 + m_2).
    \end{alignedat}
\end{equation}
取 $0! = 1$，空和取零。例如在式\eqref{eq:cg-one-one}的中心格，$j_1 = j_2 = j = 1$，三个投影全为零，允许 $k = 0,1$。两个分母都是 $1$，故两项恰好抵消。

有些角动量表按三列排列量子数，使用 Wigner $3j$ 符号。本书的系数与它的换算为
\begin{equation}
    C^{jm}_{j_1m_1,j_2m_2}
      = (-1)^{j_1-j_2+m}\sqrt{2j + 1}
       \begin{pmatrix} j_1 & j_2 & j \\ m_1 & m_2 & -m \end{pmatrix}.
\end{equation}
右侧第二行的投影和为零，故末项写 $-m$。换算约定及有限和可查
\href{https://dlmf.nist.gov/34.1}{NIST 数学函数手册第 34.1 节}和\href{https://dlmf.nist.gov/34.2}{第 34.2 节}，对称性见\href{https://dlmf.nist.gov/34.3}{第 34.3 节}。
